% Рубежный контроль №1 «Моделирование» — Вариант 2
% Источники: mathmod.pdf, Modelirovanie_1_-3.pdf, 23_03_2020_..._Моделирование_2.doc
\documentclass[12pt,a4paper]{article}
\usepackage[T2A]{fontenc}
\usepackage[utf8]{inputenc}
\usepackage[russian]{babel}
\usepackage{amsmath,amssymb}
\usepackage{geometry}
\geometry{margin=2.2cm}
\usepackage{enumitem}
\setlist{nosep,leftmargin=*}

\title{Рубежный контроль №1 «Моделирование»\\Вариант 2 — ответы}
\author{}
\date{}

\begin{document}
\maketitle

\begin{enumerate}[label=\textbf{\arabic*.}]

\item \textbf{Имитационная модель (определение).}

Имитационная модель — численный метод вычислительного эксперимента, позволяющий имитировать поведение реальной системы на ЦВМ в заданный или формируемый период времени при заданных входных данных. Для многокомпонентных систем используют дискретное модельное время и схемы квазипараллелизма/синхронизации.

\item \textbf{Феноменологическая модель (определение).}

Феноменологические модели выделяются отдельно как модели, \emph{построенные на результатах конкретного (уже оконченного) эксперимента} — то есть описывают наблюдаемую закономерность без вывода из физического механизма, и применимы в пределах условий, при которых получены экспериментальные данные.

\item \textbf{Натурный эксперимент (определение).}

Натурный эксперимент — эксперимент с реальными объектами и процессами (объект исследования в естественных условиях), в отличие от вычислительного эксперимента на модели; как правило, он ресурсоёмок и/или ограничен по безопасности и повторяемости.

\item \textbf{Классификация сложных динамических систем (перечислить).}

В классификации курса для динамических систем выделяют, в частности:
\begin{enumerate}[label=\arabic*)]
\item \textbf{простые динамические системы};
\item \textbf{структурно-сложные динамические системы};
\item \textbf{сложные динамические системы, меняющие поведение во времени};
\item \textbf{структурно-сложные гибридные системы} (сочетают структурную сложность и гибридность/разрывы поведения).
\end{enumerate}

\item \textbf{Выбор модельного времени с переменным шагом.}

Модельное время — дискретное время в единицах реального времени, обеспечивающее внутреннюю и внешнюю синхронизацию компонент модели. При выборе \textbf{переменного шага} кванты модельного времени зависят от событийной структуры процесса:
\begin{itemize}
\item если два и более событий должны быть описаны в один и тот же момент модельного времени, выполняют последовательное описание компонент с последующей коррекцией кванта;
\item если в реальной системе длительное время не возникает новых событий, шаг увеличивают (иначе имитация «простаивает»).
\end{itemize}
В практической реализации это приводит к событийной дискретизации (шаг равен времени до ближайшего события) или к адаптивной схеме изменения шага.

\item \textbf{Методика построения аналитической динамической модели.}

Методика (по курсу):
\begin{enumerate}[label=\arabic*)]
\item оценивается постановка задачи и процесс изменения объекта во времени иллюстрируется графической схемой (в моменты $t_0$, $t^{*}$, $T$) или диаграммой;
\item по схеме уточняются входные/выходные данные, выделяются промежуточные переменные;
\item определяются физико-химические воздействия, влияющие на изменение выходных данных;
\item оценивается возможность упрощения (часть воздействий не учитывается);
\item формализуется модель: составляется уравнение/система, связывающая входы и выходы (с согласованием размерностей);
\item конкретизируется описание каждого отдельного воздействия;
\item повторно оценивается упрощение: число неизвестных не должно превышать числа уравнений, уточняются коэффициенты/константы и начальные/краевые условия.
\end{enumerate}

\item \textbf{Функциональная зависимость (определение).}

Функциональная зависимость — соответствие, при котором каждой величине $x\in X$ ставится в соответствие единственное значение $y\in Y$ (в отличие от стохастической зависимости, где одному $x$ могут соответствовать различные $y$).

\item \textbf{9–10. Имеются результаты ЕГЭ 10 абитуриентов. Оценить силу связи между показателями посредством рангового коэффициента Спирмена.}

Для двух показателей $X$ и $Y$ ранговый коэффициент Спирмена вычисляют как корреляцию Пирсона между рангами:
\[
\rho_s=\mathrm{corr}(\mathrm{rank}(X),\mathrm{rank}(Y)).
\]
При наличии совпадающих значений используются усреднённые ранги.

\textbf{Связь между РЯ и М.}

\textbf{Шаг 1. Ранги (по возрастанию баллов; при совпадениях — средний ранг).}

Для РЯ (82, 78, 85, 96, 98, 56, 84, 88, 72, 66) получаем ранги:
\[
r_{\text{РЯ}}=(5,\ 4,\ 7,\ 9,\ 10,\ 1,\ 6,\ 8,\ 3,\ 2).
\]

Для М (76, 82, 90, 98, 96, 76, 78, 98, 80, 82) ранги с учётом повторов:
\[
r_{\text{М}}=(1.5,\ 5.5,\ 7,\ 9.5,\ 8,\ 1.5,\ 3,\ 9.5,\ 4,\ 5.5).
\]

\textbf{Шаг 2. Корреляция Пирсона для рангов.}
Обозначим $\bar r_X=\frac1n\sum r_{X,i}$, $\bar r_Y=\frac1n\sum r_{Y,i}$ (здесь $n=10$, поэтому $\bar r_X=\bar r_Y=5.5$).
Тогда
\[
\rho_s=\frac{\sum_{i=1}^n (r_{X,i}-\bar r_X)(r_{Y,i}-\bar r_Y)}
{\sqrt{\sum_{i=1}^n (r_{X,i}-\bar r_X)^2}\ \sqrt{\sum_{i=1}^n (r_{Y,i}-\bar r_Y)^2}}.
\]
Вычисления дают:
\[
\sum (r_{\text{РЯ}}-5.5)(r_{\text{М}}-5.5)=60.0,\quad
\sum (r_{\text{РЯ}}-5.5)^2=82.5,\quad
\sum (r_{\text{М}}-5.5)^2=81.0,
\]
откуда
\[
\rho_s(\text{РЯ},\text{М})=\frac{60.0}{\sqrt{82.5\cdot 81.0}}\approx 0.734.
\]

\textbf{Ответ:} $\rho_s(\text{РЯ},\text{М})\approx 0.734$.
\[
\rho_s(\text{РЯ},\text{М})\approx 0.734.
\]

\textbf{Связь между Ф и И.}

\textbf{Шаг 1. Ранги.}

Для Ф (85, 80, 88, 92, 100, 74, 72, 90, 65, 69):
\[
r_{\text{Ф}}=(7,\ 5,\ 8,\ 9,\ 10,\ 4,\ 3,\ 6,\ 1,\ 2).
\]

Для И (68, 82, 94, 100, 84, 64, 70, 90, 74, 78):
\[
r_{\text{И}}=(2,\ 7,\ 9,\ 10,\ 8,\ 1,\ 3,\ 6,\ 4,\ 5).
\]

\textbf{Шаг 2. Корреляция рангов.}
\[
\sum (r_{\text{Ф}}-5.5)(r_{\text{И}}-5.5)=53.5,\quad
\sum (r_{\text{Ф}}-5.5)^2=82.5,\quad
\sum (r_{\text{И}}-5.5)^2=82.5,
\]
поэтому
\[
\rho_s(\text{Ф},\text{И})=\frac{53.5}{82.5}\approx 0.648.
\]

\textbf{Ответ:} $\rho_s(\text{Ф},\text{И})\approx 0.648$.
\[
\rho_s(\text{Ф},\text{И})\approx 0.648.
\]

\textbf{Связь между остальными парами предметов (аналогично, по рангам).}

\begin{itemize}
\item \textbf{РЯ и Ф.} Ранги $r_{\text{РЯ}}=(5,4,7,9,10,1,6,8,3,2)$, $r_{\text{Ф}}=(6,5,7,9,10,4,3,8,1,2)$.
\[
\sum (r_{\text{РЯ}}-5.5)(r_{\text{Ф}}-5.5)=70.5,\quad
\sum (r_{\text{РЯ}}-5.5)^2=82.5,\quad
\sum (r_{\text{Ф}}-5.5)^2=82.5,
\]
\[
\rho_s(\text{РЯ},\text{Ф})=\frac{70.5}{82.5}\approx 0.855.
\]
\item \textbf{РЯ и И.} Ранги $r_{\text{РЯ}}=(5,4,7,9,10,1,6,8,3,2)$, $r_{\text{И}}=(2,6,9,10,7,1,3,8,4,5)$.
\[
\sum (r_{\text{РЯ}}-5.5)(r_{\text{И}}-5.5)=59.5,\quad
\sum (r_{\text{РЯ}}-5.5)^2=82.5,\quad
\sum (r_{\text{И}}-5.5)^2=82.5,
\]
\[
\rho_s(\text{РЯ},\text{И})=\frac{59.5}{82.5}\approx 0.721.
\]
\item \textbf{М и Ф.} Ранги $r_{\text{М}}=(1.5,5.5,7,9.5,8,1.5,3,9.5,4,5.5)$, $r_{\text{Ф}}=(6,5,7,9,10,4,3,8,1,2)$.
\[
\sum (r_{\text{М}}-5.5)(r_{\text{Ф}}-5.5)=54.5,\quad
\sum (r_{\text{М}}-5.5)^2=81.0,\quad
\sum (r_{\text{Ф}}-5.5)^2=82.5,
\]
\[
\rho_s(\text{М},\text{Ф})=\frac{54.5}{\sqrt{81.0\cdot 82.5}}\approx 0.667.
\]
\item \textbf{М и И.} Ранги $r_{\text{М}}=(1.5,5.5,7,9.5,8,1.5,3,9.5,4,5.5)$, $r_{\text{И}}=(2,6,9,10,7,1,3,8,4,5)$.
\[
\sum (r_{\text{М}}-5.5)(r_{\text{И}}-5.5)=77.5,\quad
\sum (r_{\text{М}}-5.5)^2=81.0,\quad
\sum (r_{\text{И}}-5.5)^2=82.5,
\]
\[
\rho_s(\text{М},\text{И})=\frac{77.5}{\sqrt{81.0\cdot 82.5}}\approx 0.948.
\]
\end{itemize}

\textbf{Итог по всем 6 парам.}
\[
\rho_s(\text{РЯ},\text{М})\approx 0.734,\quad
\rho_s(\text{РЯ},\text{Ф})\approx 0.855,\quad
\rho_s(\text{РЯ},\text{И})\approx 0.721,
\]
\[
\rho_s(\text{М},\text{Ф})\approx 0.667,\quad
\rho_s(\text{М},\text{И})\approx 0.948,\quad
\rho_s(\text{Ф},\text{И})\approx 0.648.
\]
\textbf{Вывод:} во всех парах связь \emph{положительная}. Самая сильная — между \textbf{М} и \textbf{И} ($\approx 0.948$), затем \textbf{РЯ} и \textbf{Ф} ($\approx 0.855$). Самая слабая из шести — между \textbf{Ф} и \textbf{И} ($\approx 0.648$) и между \textbf{М} и \textbf{Ф} ($\approx 0.667$).

\end{enumerate}
\end{document}

